%% ma: L12-B17-0019
%% lop: 12
%% chuong: 5
%% bai: 17
%% dang: 12.17.13
%% loai: TLN
%% muc: VDC
%% dap_an: "14"
%% nguon: "(NBV)-10 ĐỀ KIỂM TRA HK2-TOÁN 12-DUNG CHUNG 3 SÁCH.pdf (H02, Đề số 2), Phần 3 Câu 2"
%% kien_thuc: "Bài 17 (mặt cầu), Bài 6, 8 (vectơ trong không gian)"
%% trang_thai: cho-duyet
%% ly_do: "Dạng mới 12.17.13 chờ duyệt."
%% ghi_chu: "Đáp án gốc 14 khớp (sympy)."
\begin{ex}
Trong không gian $Oxyz$, cho ba điểm $A(8;5;-11)$, $B(5;3;-4)$, $C(1;2;-6)$ và mặt cầu $(S)$ có phương trình
\[(x-2)^2+(y-4)^2+(z+1)^2=36.\]
Gọi $M(a;b;c)$ là điểm trên $(S)$ sao cho
\[\left|\overrightarrow{MA}-\overrightarrow{MB}-\overrightarrow{MC}\right|\]
đạt giá trị lớn nhất. Hãy tìm $a+b$.
\shortans{14}
\loigiai{Mặt cầu $(S)$ có tâm $I(2;4;-1)$, bán kính $R=6$. Gọi $N$ là điểm thỏa $\overrightarrow{NA}-\overrightarrow{NB}-\overrightarrow{NC}=\vec0$, suy ra $N(-2;0;1)$.

Khi đó
\[\overrightarrow{MA}-\overrightarrow{MB}-\overrightarrow{MC}=\left(\overrightarrow{NA}-\overrightarrow{NB}-\overrightarrow{NC}\right)-\overrightarrow{MN}=\overrightarrow{NM},\]
nên $\left|\overrightarrow{MA}-\overrightarrow{MB}-\overrightarrow{MC}\right|=MN$.

Vì $IN=6=R$ nên $N\in(S)$, do đó $MN\le2R=12$, dấu bằng khi $M$ đối xứng với $N$ qua $I$, tức $M(6;8;-3)$. Vậy $a+b=14$.}
\end{ex}
